\paragraph{A member's balances can depend on state held by the group.}
The corporate group introduces a different reason to preserve the native
representation. Its parent and subsidiary have separate balances, but their
state also includes shared inventory lots, returns, a compound instrument,
tax calculations and ownership changes. Copying each balance dictionary into
an independent company book would leave those dependencies behind. A later
event could then use a valid-looking opening balance without the facts needed
to determine its accounting consequences.

The reference adapter therefore selects a member within the complete group
state. A selector names the exact state record and the path
\code{/books/P} or \code{/books/S}; its identity differs from the identity of
the group record itself. The common company's reference points to that
selector. The result retains the selected balances, the whole native state
and its information view together. Record schemas, dependency references and
source bytes are checked before a selection is returned. The selector gives
access to a member's balances without asserting that they form a separately
executable book.

The two retained snapshots also show why time belongs in the relationship.
On 1 January 2025, the supplied parent ownership fraction is four-fifths and
the parent and subsidiary cash balances are 500 and 100. On 31 December 2025,
ownership is nine-tenths and cash is 491.5 and 152.5. Binding either fraction
to the entire supplied control interval would erase an ownership change.
The adapter instead attaches it to the selected snapshot day. It preserves
the supplied control premise and the archival, synthetic information view;
it does not determine control, recalculate consolidation or admit current
source evidence.

\paragraph{An elimination worksheet leaves the legal books intact.}
Separately valid company books permit a different operation. Consider a new
constructed example in USD units at 31 December 2025. Parent P holds cash 80,
an investment in subsidiary S of 100 and a receivable from S of 20; its equity
is 200. S holds cash 70 and equipment 50, owes P 20 and has capital of 100.
Both books balance. Assume that P controls and fully owns S, the books use
compatible accounting policies and measurement bases, the reciprocal balances
refer to the same obligation, and the investment equals the subsidiary's
capital with no acquisition difference. These are supplied facts of the
illustration. They do not describe the earlier four-fifths or nine-tenths
ownership case.

Adding the legal books gives assets of 320, liabilities of 20 and equity of
300. The worksheet then records two explicit adjustments. First, it credits
the investment by 100 and debits the subsidiary capital by 100. Second, it
credits the receivable by 20 and debits the payable by 20. Under the declared
signed-balance convention, each adjustment sums to zero. The resulting assets
are cash 150 and equipment 50; liabilities are zero and equity is 200. Cash
does not change. Neither legal book receives an entry: the adjustments belong
to the separate worksheet, alongside the unadjusted contributions.

The implementation requires a complete map from each local account to the
declared common chart. It checks the legal owners, case, currency, period,
snapshot date and classifications, and preserves each book's memory and cash
flow history. Every paired adjustment names its source accounts, expected
full balances and economic premise. A source account can be used only once;
the pair must come from different books and contain exactly opposite,
nonzero balances. This first operation supports the two declared pair roles,
receivable--payable and investment--equity. It refuses cash elimination,
partial allocations, unmatched balances, unresolved components and unsupported
sector books rather than supplying a balancing amount. A worksheet with no
declared adjustments remains an unadjusted combination.

Two versions of the illustration use different local account names and the
same explicit common-chart meanings. Both reproduce the complete signed
unadjusted balances, elimination deltas and adjusted balances; replay also
checks that the legal books remain unchanged. These checks establish the
finite worksheet calculation. The truth of the supplied economic premises
still requires evidence, and classification equality does not prove policy
equivalence. Noncontrolling interests, goodwill, tax, foreign exchange,
intragroup profit and period cash-flow eliminations require additional methods.
The result therefore remains a snapshot worksheet, not a complete consolidated
statement package or an executable group book.

The reference selections and worksheet cases are retained in two directories
under \path{data/accounting-master/}:
\begin{flushleft}\small
\path{native-group-binding-20261025}\\
\path{compatible-book-composition-20261025}
\end{flushleft}

